\section{Analysis \& Discussion}
\label{sec:discussion}

This section interprets the comparative results across the three scanner
categories, identifies the key patterns that emerge from the evaluation, and
acknowledges the limitations and threats to validity that constrain
the conclusions we can draw.

%% ----------------------------------------------------------------
\subsection{Key Findings}
\label{sec:key-findings}

\paragraph{A clear three-tier hierarchy emerges.}
The most striking result is the emergence of three distinct performance
tiers.  The security-specialized scanner (Kolega, F3\,=\,73.0 strict)
substantially outperforms even the best general-purpose LLM (Sonnet~4.6,
F3\,=\,51.7), which in turn achieves nearly 3$\times$ the F3 score of
the best rule-based tool (Semgrep, F3\,=\,17.7).

This hierarchy suggests that general-purpose LLM reasoning, while a
major advance over pattern matching, is not sufficient for
production-grade vulnerability detection.  Purpose-built security
architecture---combining LLM reasoning with domain-specific
design---delivers meaningfully higher recall.

We note two important caveats. First, the security-specialized category currently contains only one system (Kolega); additional tools in this category would be needed to confirm a distinct performance tier rather than a single outlier. Second, the general-purpose LLM tier exhibits high internal variance, with F3 scores ranging from 20.6 (Grok~3) to 51.7 (Sonnet~4.6) --- a spread larger than the gap between the best GP-LLM and the best rule-based tool. Some GP-LLM scanners (Grok~3, Haiku single-turn) perform at or below rule-based SAST levels, suggesting that model choice and configuration matter as much as the LLM-vs-SAST distinction.

The per-CWE breakdown reveals where these advantages originate.
General-purpose LLMs dramatically outperform rule-based SAST on
vulnerability classes that require semantic understanding: SQL injection
(91\% vs.\ 32\% recall), command injection (78\% vs.\ 24\%), and
insecure deserialization.  Rule-based SAST tools remain competitive only
on issues that reduce to syntactic patterns (e.g., hardcoded secrets),
but even on those categories their overall recall remains low.

\paragraph{Security-specialized architecture closes the gap that general-purpose LLMs cannot.}
The 41\% F3 gap between Kolega (73.0) and the best general-purpose LLM (Sonnet~4.6, 51.7) warrants examination. General-purpose LLMs bring powerful code reasoning but lack security-specific optimizations: they receive a generic prompt, process code through a general-purpose context window, and produce findings without domain-specialized validation. Security-specialized systems can incorporate vulnerability-aware code analysis pipelines, security-specific retrieval and context assembly, and domain-optimized detection strategies that compound to produce meaningfully higher recall. Critically, Kolega achieves this while completing 26 of 26 repositories with zero timeouts --- demonstrating that purpose-built architecture solves the reliability problems that plague frontier general-purpose models (Opus: 36\% timeout rate, Minimax: 31\%). The combination of higher recall \emph{and} higher reliability is the defining advantage of the security-specialized category.

\paragraph{Where Kolega falls short.}
Despite leading the aggregate rankings, Kolega's precision of 0.388 is the lowest of any evaluated scanner --- meaning approximately 61\% of its findings are false positives. This is a substantial triage burden: for every true vulnerability found, an analyst must review roughly 1.6 additional false alarms. Teams with limited analyst bandwidth may reasonably prefer a higher-precision scanner such as Sonnet~4.6 (precision 0.785) at the cost of lower recall. Kolega also underperforms general-purpose LLMs on several individual repositories: on \texttt{realvuln-python-app} (Kolega F3\,=\,55.8 vs.\ Opus~4.6 at 83.5), \texttt{realvuln-vulnpy} (Kolega 62.6 vs.\ Gemini~3.1~Pro at 86.9), and \texttt{realvuln-damn-vulnerable-flask-application} (Kolega 67.0 vs.\ Opus~4.6 at 77.4). These repos suggest that Kolega's advantage is not universal --- on well-structured codebases with canonical vulnerability patterns, general-purpose LLMs can match or exceed it.

\paragraph{Agentic tool use is highly cost-effective.}
Comparing agentic and single-turn modes of the same underlying model
provides a controlled test of the value of tool use in vulnerability
detection.  Agentic Haiku achieves an F3 of 37.1 (strict) versus 25.2 for
single-turn Haiku---a 47\% improvement---while incurring only 6\% higher
cost.  The ability to read files, grep for patterns, and iteratively
explore a codebase allows the model to gather context that a single
prompt cannot provide.  This is consistent with the broader finding that
retrieval-augmented and tool-augmented LLM workflows outperform
prompt-only setups in code understanding tasks.  From a practical
standpoint, the marginal cost of agentic execution is small relative to
the gain in detection quality, making it the dominant strategy for
LLM-based scanning.

\paragraph{Bigger models do not always win; reliability matters.}
One might expect a monotonic relationship between model capability
(and cost) and benchmark performance.  The data tell a more nuanced
story.  Opus~4.6 achieves strong per-repo scores on the repositories
it completes, but its 36\% timeout rate (completing only 19 of 26
repos) drags its strict F3 down to 47.5---below Sonnet~4.6's 51.7,
which reflects fuller coverage.  A similar pattern emerges with
Minimax, which times out on 31\% of repos.  Conversely,
Kimi~K2.5 and Minimax~M2.7 punch above their weight on
cost-efficiency metrics, delivering competitive F3 scores at
substantially lower per-repo cost.  The implication for practitioners is
that scanner selection should optimize for \emph{reliability-adjusted}
performance rather than peak capability: a model that reliably completes
every scan is more valuable than one that occasionally excels but
frequently fails to return results.

\paragraph{The precision--recall tradeoff is stark and F3 captures it correctly.}
Grok~4.20 Reasoning exemplifies the conservative end of the spectrum:
its precision of 0.927 is the highest in the benchmark, but its recall
of 0.263 means it misses nearly three-quarters of true vulnerabilities.
Most LLM scanners cluster in the 0.6--0.8 precision range with
substantially higher recall, which the F3 metric---designed to weight
recall nine times more heavily than precision---correctly rewards.  Among
rule-based tools, SonarQube illustrates the opposite failure mode:
respectable precision (0.611) paired with negligible recall (0.065),
yielding an F3 of just 6.8.  The F3 metric aligns well with
practitioner priorities in security scanning, where a single missed
vulnerability can lead to a catastrophic breach while false positives
merely cost analyst time.  Benchmarks that report
only precision or only detection rate would paint a misleading picture
of scanner utility.

%% ----------------------------------------------------------------
\subsection{Limitations}
\label{sec:limitations}

We identify seven limitations that scope the conclusions of this work:

\begin{itemize}
  \item \textbf{Python only.}
    Version~1.0 of \textsc{RealVuln} covers only Python repositories.
    Scanner performance may differ substantially on statically typed
    languages (Java, Go) or memory-unsafe languages (C/C++) where
    vulnerability classes such as buffer overflows are prevalent.
    Extension to additional languages is planned for v2
    (Section~\ref{sec:living-benchmark}).

  \item \textbf{Type~1 repositories only.}
    All 26 repos are intentionally vulnerable applications with
    artificially high vulnerability density.  Production codebases have
    sparser vulnerabilities embedded in larger, more complex code, which
    may alter both precision and recall characteristics.  Type~2--5
    targets (Section~\ref{sec:target-types}) will address this gap.

  \item \textbf{Vendor conflict of interest.}
    Kolega, which leads the benchmark at F3\,=\,73.0 (strict), is also the
    organization publishing this work.  We mitigate this concern by
    open-sourcing the entire benchmark: ground truth labels, scoring
    code, scanner outputs, and the dashboard.  Any researcher can
    independently verify our results or challenge the ground truth.

  \item \textbf{Ground truth subjectivity.}
    Hand-labeling vulnerabilities inherently involves judgment calls,
    particularly for borderline cases such as informational findings or
    context-dependent misconfigurations.  We publish the labeling
    evidence and rationale for every ground truth entry and invite
    community challenges via the project's issue tracker.

  \item \textbf{LLM non-determinism.}
    LLM-based scanners produce variable output across runs.  We report
    multi-run statistics (mean, standard deviation) where available, but
    cost constraints prevented multi-run evaluation for every model
    configuration.  Single-run results should be interpreted with
    appropriate uncertainty.

  \item \textbf{Default scanner configurations.}
    All scanners are evaluated using default or vendor-recommended
    settings.  Tuned configurations---custom rules for Semgrep, adjusted
    quality profiles for SonarQube, or modified system prompts for
    LLM scanners---could yield different results.  We report the exact
    configuration used for each scanner to enable reproduction.

  \item \textbf{Timeouts and incomplete scans.}
    Opus~4.6 timed out on approximately 36\% of repositories and
    Minimax~M2.7 on approximately 31\%.  These incomplete scans reduce
    effective coverage and are reflected in the strict micro F3 metric,
    which treats unscanned repos as zero.  The gap between lenient and
    strict scores quantifies the reliability penalty.

  \item \textbf{Missing scanner families.}
    OpenAI models (GPT-4o, o3) are absent from this evaluation due to API restrictions on automated security scanning at the time of evaluation. Their inclusion would broaden the general-purpose LLM tier. Similarly, open-weight models (Llama, Mistral) are not evaluated. We plan to include both in v2.

  \item \textbf{Single-turn vs.\ agentic comparison is limited.}
    Only Claude Haiku~4.5 was evaluated in both single-turn and agentic modes. The 47\% F3 improvement we report is based on this single model pair; different models may benefit more or less from agentic scaffolding.
\end{itemize}

%% ----------------------------------------------------------------
\subsection{Threats to Validity}
\label{sec:threats}

\paragraph{Data contamination.}
The 26 target repositories are publicly available on GitHub, and it is
likely that some or all of them appeared in the training data of the LLM
scanners we evaluate.  If a model has memorized known vulnerabilities in
these repos, its recall could be inflated relative to performance on
unseen code.  We mitigate this threat in two ways.  First, the 120
false-positive traps in the ground truth test whether scanners
\emph{discriminate} between vulnerable and safe code; memorizing that a
file contains vulnerabilities does not help a scanner avoid flagging a
non-vulnerable code path in the same file.  Second, the v2 roadmap
(Section~\ref{sec:living-benchmark}) includes repositories published
after model training cutoffs, which will provide a contamination-free
evaluation set.

\paragraph{Selection bias.}
Twenty-six repositories, while substantial for a hand-labeled benchmark,
do not exhaustively represent the Python vulnerability landscape.
Certain CWE families (e.g., cryptographic weaknesses, race conditions)
are underrepresented in the current corpus.  The living-benchmark
roadmap addresses this through planned expansion to additional repos
and target types, aiming for broader CWE and application-domain
coverage in future versions.
